\chapter{Technologies, frameworks et bibliothèques}
\minitoc

Ce chapitre s'adresse à un lecteur technique ou à un évaluateur. Il répond à
trois questions : sur quoi la plateforme est bâtie, quelles bibliothèques elle
emploie et où, et quelle est la nature exacte de l'« intelligence » qu'elle met
en œuvre.

\section{Le socle}

\begin{table}[h]
\centering
\begin{tabular}{@{}p{3.6cm}p{2.4cm}p{7.8cm}@{}}
\toprule
\textbf{Élément} & \textbf{Version} & \textbf{Rôle} \\
\midrule
Python & 3.11+ & Langage unique du projet. \\[5pt]
FastAPI & 0.141 & Cadre applicatif web : les 91 routes de l'API. \\[5pt]
Pydantic & 2.13 & Validation des charges utiles entrantes. \\[5pt]
Uvicorn & 0.52 & Serveur d'application. \\[5pt]
SQLite & intégré & Base de données, en mode journal WAL. \\[5pt]
PyYAML & 6.0 & Lecture des playbooks. \\[5pt]
ReportLab & 5.0 & Rendu PDF des rapports. \\[5pt]
python-docx & 1.2 & Rendu Word des rapports. \\[5pt]
HTML/CSS/JS & --- & Console web, sans aucun cadre applicatif. \\
\bottomrule
\end{tabular}
\caption{Le socle technique}
\end{table}

\begin{tcolorbox}[colback=green!4!white, colframe=darkgreen,
  title=\textbf{Bon à savoir --- pourquoi si peu de dépendances},
  fonttitle=\bfseries, boxrule=0.6pt]
ReportLab et python-docx sont écrits en Python pur : aucune bibliothèque
système à installer, et la génération des rapports reste possible sur un site
coupé d'Internet. La console est écrite en JavaScript sans cadre applicatif ni
étape de construction — elle se sert telle quelle. Ce sont des contraintes du
Centre, pas des préférences.
\end{tcolorbox}

\section{Où chaque bibliothèque intervient}

Le relevé qui suit porte sur les imports réels du dépôt, module par module.

\begin{table}[h]
\centering
\begin{tabular}{@{}p{4.4cm}p{3.6cm}p{5.8cm}@{}}
\toprule
\textbf{Module} & \textbf{Tierces} & \textbf{Standard Python employé} \\
\midrule
\texttt{ingestion/} & \textbf{aucune} & \texttt{re}, \texttt{datetime} \\[4pt]
\texttt{domain/} & \textbf{aucune} & \texttt{dataclasses}, \texttt{enum},
\texttt{uuid} \\[4pt]
\texttt{enrichment/} & \textbf{aucune} & \texttt{math}, \texttt{collections} \\[4pt]
\texttt{orchestration/} & PyYAML & \texttt{hashlib}, \texttt{ipaddress},
\texttt{re} \\[4pt]
\texttt{actuators/} & \textbf{aucune} & \texttt{abc}, \texttt{uuid} \\[4pt]
\texttt{detection/} & \textbf{aucune} & \texttt{math}, \texttt{json} \\[4pt]
\texttt{audit/} & \textbf{aucune} & \texttt{hashlib}, \texttt{sqlite3},
\texttt{threading} \\[4pt]
\texttt{persistence/} & \textbf{aucune} & \texttt{sqlite3}, \texttt{json} \\[4pt]
\texttt{api/} & FastAPI, Pydantic & frontière HTTP seulement \\[4pt]
\texttt{reporting/} & ReportLab, python-docx & sortie documentaire seulement \\[4pt]
\texttt{llm/} & \texttt{anthropic} & \textbf{facultatif, non installé par
défaut} \\
\bottomrule
\end{tabular}
\caption{Bibliothèques par module}
\end{table}

Ce tableau porte le résultat le plus notable du projet : \textbf{toute la
chaîne de décision} — de la normalisation d'une observation à son inscription
au registre — \textbf{ne dépend que de la bibliothèque standard}, à l'exception
de PyYAML qui ne fait que lire des fichiers de playbook.

Trois conséquences :

\begin{itemize}[leftmargin=*]
  \item le comportement est \textbf{entièrement lisible dans le dépôt} : aucun
        moteur de règles tiers, aucun ORM, aucune boîte noire~;
  \item la plateforme tourne sur un site coupé d'Internet~;
  \item la surface d'attaque de la chaîne critique est celle de Python
        lui-même.
\end{itemize}

\section{Les algorithmes mis en œuvre}

\begin{description}[leftmargin=3.4cm, style=nextline]
  \item[BM25 Okapi] Recherche documentaire pour l'ancrage
        ($k_1 = 1{,}5$, $b = 0{,}75$), sur index inversé. Écrit avec le module
        \texttt{math}. C'est ce calcul qui autorise ou refuse l'action.
  \item[Welford] Calcul incrémental de moyenne et de variance pour les profils
        comportementaux. Permet de tenir une référence glissante sans conserver
        l'historique.
  \item[SHA-256] Empreintes de déduplication des observations, chaînage du
        journal d'audit, signature des menaces inédites, empreinte de la
        politique compilée.
  \item[PBKDF2-HMAC-SHA256] Dérivation des mots de passe, 600\,000 itérations.
  \item[Score de risque] Somme pondérée de six composantes explicables, sans
        aucun poids appris.
\end{description}

\section{L'intelligence : de quelle nature ?}

C'est la question qui vient toujours. La réponse doit être précise, parce que
trois notions différentes sont souvent confondues : \emph{autonome},
\emph{intelligence artificielle}, et \emph{modèle de langage}.

\subsection{Ce n'est pas un agent fondé sur un modèle de langage}

La plateforme n'est pas une boucle « instruction, outil, observation » pilotée
par un modèle. C'est un \textbf{pipeline déterministe} de neuf étapes.
\textbf{Aucun modèle de langage ne décide quoi que ce soit.}

\subsection{Trois familles portent l'intelligence}

\begin{description}[leftmargin=*, style=nextline]
  \item[Apprentissage statistique non supervisé] \texttt{detection/ueba/}.
        Profil comportemental par entité : moyenne et écart-type glissants,
        scores d'écart. Le scoreur \textbf{s'abstient sous sept observations},
        plutôt que de qualifier d'anormal ce qu'il n'a jamais vu. Le système
        construit bien son modèle du normal à partir des données : c'est de
        l'apprentissage.
  \item[Recherche d'information] \texttt{enrichment/}. BM25 sur un fonds de
        vingt-huit fiches de référence. Ce n'est pas de l'apprentissage, mais
        c'est ce qui commande l'autorisation d'agir.
  \item[Raisonnement symbolique] \texttt{domain/taxonomy.py},
        \texttt{orchestration/}. Base de connaissances explicite — vingt-deux
        types d'attaques, trente-huit gestes catalogués, playbooks, table des
        objectifs de confinement — et moteur d'inférence. C'est la branche
        historique de l'intelligence artificielle, celle des systèmes experts,
        et c'est elle qui porte l'essentiel des décisions.
\end{description}

\subsection{Où un modèle de langage peut intervenir}

Deux emplacements, tous deux facultatifs et hors du chemin d'exécution :

\begin{enumerate}[leftmargin=*]
  \item \textbf{la mise en forme des réponses de l'assistant} — les chiffres
        restent ceux calculés par la plateforme, le modèle ne fait que
        rédiger~;
  \item \textbf{le conseil après un refus} — le modèle peut rédiger la
        justification, voire désigner l'alternative, mais
        \textbf{exclusivement parmi les gestes déjà calculés} par la recherche
        de substitution.
\end{enumerate}

\begin{tcolorbox}[colback=green!4!white, colframe=darkgreen,
  title=\textbf{Bon à savoir --- l'invariant qui rend cela acceptable},
  fonttitle=\bfseries, boxrule=0.6pt]
Le modèle \textbf{choisit dans une liste, il ne la fabrique pas}. Tout candidat
a déjà passé les trois conditions de la substitution : exécutable seul,
entièrement réversible, strictement moins engageant que le geste écarté. Un
modèle qui répondrait n'importe quoi ne peut pas faire sortir un geste
dangereux — sa réponse est confrontée à la liste et rejetée si elle n'y figure
pas. Un test le vérifie avec un fournisseur qui propose délibérément d'effacer
un disque.
\end{tcolorbox}

Par défaut, le fournisseur configuré est \texttt{offline} et le paquet
\texttt{anthropic} n'est pas installé : \textbf{aucun appel à un modèle n'a
lieu}. Pour l'activer :

\begin{lstlisting}[language=bash]
$ pip install anthropic
$ export CIRT_LLM_PROVIDER=anthropic
$ export ANTHROPIC_API_KEY=...
\end{lstlisting}

\subsection{Pourquoi ce choix se défend}

Sans validation humaine en amont, la seule défense d'une action est sa
justification \emph{a posteriori}. Un playbook rend une réponse exacte : telle
règle, de tel fichier, dans telle version. Un modèle génératif rend au mieux
une reconstruction plausible de son propre raisonnement — ce qui n'est pas la
même chose, et ne se vérifie pas.

S'y ajoutent le déterminisme — la même menace doit produire la même réponse —
et l'énumérabilité du périmètre d'action : on doit pouvoir dire exactement ce
que le système est capable de faire.

\section{Architecture logicielle}

L'architecture est hexagonale. La couche \texttt{domain/} — modèle métier,
invariants, taxonomie — \textbf{ne dépend de rien} : ni de la base, ni du web,
ni d'aucune bibliothèque tierce. Les couches extérieures dépendent d'elle, et
jamais l'inverse.

C'est ce qui permet de remplacer SQLite, FastAPI ou un actuateur sans toucher
aux règles métier — et c'est aussi ce qui rend la suite de tests rapide : la
plus grande part s'exécute sans base et sans serveur.

\section{Qualité}

\begin{itemize}[leftmargin=*]
  \item \textbf{561 tests} — unitaires, d'intégration, et de recette (les
        critères CR du cahier des charges).
  \item \textbf{Ruff} pour le style et les erreurs courantes, sur le code
        source comme sur les tests.
  \item Des \textbf{tests de garde} qui figent les invariants : absence de
        route de validation préalable dans le schéma de l'API, absence de
        dialogue natif dans l'interface, traduction obligatoire de tout geste
        du catalogue.
\end{itemize}
